\chapter{Conclusion}
\label{ch:conclusion}

\section{Answers to the Research Questions}

\paragraph{RQ1.} A frequent batch auction engine is fully specified by
five design blocks: the batch interval, the disclosure regime during
accumulation, the order flow segregation architecture, the clearing and
rationing rule, and settlement and residual handling. These are not
optional extras layered onto the mechanism; they are decisions that a
continuous engine also makes but makes implicitly, and that become
visible and contestable only once serial processing is abandoned.
\Cref{ch:design} enumerated the options within each block and derived
the behaviour each induces, consolidating the result in the trade-off
matrix of \Cref{tab:tradeoff}.

Three findings stand out from that analysis. First, the only unambiguous
structural gain in the space is the combination of discrete accumulation
and uniform pricing; every other choice trades one distortion for
another, and the design question is which distortion the venue prefers
rather than which option is undistorted. Second, the blocks interact, so
they cannot be optimised independently: disclosure determines how much
pressure falls on the interval boundary, interval length interacts with
instrument liquidity, and the residual policy determines whether an
order can accrue persistence across batches. Third, a venue can
implement frequent batch auctions exactly as proposed and still have a
speed race, because price-time rationing inside the batch reintroduces
one through the allocation rule. This last point is the sharpest
practical implication of the chapter.

Two questions in the space have no settled answer and were developed
here as contributions rather than summaries. The maker--taker
distinction dissolves under simultaneous uniform-price clearing, and
four candidate redefinitions were set out, of which persistence across
batches has the strongest economic claim. And the price improvement that
uniform clearing produces relative to submitted limits is not a fixed
property of the mechanism but a quantity whose magnitude and
distribution depend on the clearing and rationing rules chosen.

\paragraph{RQ2.} \Cref{ch:results} reports the paired comparison of the
two engines on identical order flow across liquidity and transaction
costs, price discovery, and execution and allocation, together with the
interval sweep and the rationing comparison.
\ph[summary of measured findings to be inserted once the simulation
results are populated]

\section{Limitations}
\label{sec:limitations}

The central limitation is intrinsic to the design rather than incidental
to its execution. A replay simulation holds order flow constant, which
is exactly what makes the comparison clean, and by the same token it
prevents participants from responding to the mechanism. The orders
replayed here were submitted under a continuous engine: their sizes were
not chosen against a pro rata rationing rule, their arrival times were
not chosen against a batch boundary, and their limit prices were not
chosen under uncertainty about a clearing price. Every equilibrium
prediction of the FBA literature --- narrower quoted spreads, size
inflation, boundary concentration, bid shading --- is therefore out of
reach, and what the simulation delivers is the mechanical baseline
against which such responses would have to be measured.

Four further limitations qualify the results. The data come from a
single venue and a limited number of sessions, so magnitudes should not
be transported to other asset classes. Several metrics require
adaptation to a batch setting, and the adapted quantity is not always
the same object as the original, most notably the quoted spread and the
pre-trade midpoint reference. The engines are per-pair, so the
multi-asset clearing problem developed in \Cref{sec:clearing} is part of
the design analysis but not of the measurement. And the dataset reflects
what one non-validating node observed, which is not necessarily what the
venue's own engine processed.

\section{Future Work}
\label{sec:futurework}

\paragraph{Strategic agent populations.} The natural extension is to
replace replayed flow with agents that respond to the mechanism. A
population containing liquidity providers who requote when their quote
becomes stale, informed traders in the sense of \citet{kyle1985continuous},
and uninformed liquidity demanders would make the equilibrium
predictions P3, P5, P6, and P7 testable
\citep{lebaron2006agent}. The obvious cost is that results would then
depend on the behavioural assumptions built into the agents rather than
on observed behaviour, which is why the two designs are complements: the
replay establishes what the mechanism does mechanically, and the agent
model establishes what it does in equilibrium under stated assumptions.
A hybrid in which historical flow supplies the background and a small
set of strategic agents is layered on top would preserve much of the
realism of the former while permitting adaptation at the margin.

\paragraph{Hybrid architectures.} The comparison here is binary, but the
design space is not. A venue could run continuous matching with periodic
batch clearings for the marginal order set, batch only during volatile
periods, or batch only some instruments. Each of these is a point
between the two poles and can be evaluated with exactly the harness
built for this thesis, since the engine is substitutable at one
interface.

\paragraph{The multi-asset case.} Implementing the multi-asset clearing
problem of \eqref{eq:multiflow} and measuring how much cross-pair
arbitrage rent coherent clearing actually removes on real flow would
connect the microstructure framing of this thesis with the decentralised
exchange literature it draws on. This is a well-defined empirical
question that the present engines cannot answer because they clear one
pair at a time.

\paragraph{Fee design under batching.} Finally, the maker--taker problem
identified in \Cref{sec:segregation} could be studied directly by
implementing each of the four candidate definitions in the engine and
measuring how the resulting fee incidence is distributed across
participant types on the same flow. This is mechanical rather than
behavioural, which means it is within reach of the design used here, and
it addresses a question every venue considering batching will have to
answer.
